% [Lynn] Introduction (rubric: 1.5 marks).
% Top band: motivates the problem, states clear research question(s), and LISTS the contributions.
% FB1 (proposal feedback): the research question must be easy to find: one sentence, set apart,
%   precise about WHAT is predicted (delivery time, not delay) and WHEN (at order placement).
% Lec 2 Stage 1: who decides what with the answer (ops team: courier allocation, promise padding),
%   success criteria (ML: >= 20% MAE cut vs median; business: lower late rate at the same mean promise).
\section{Introduction}

On a food delivery platform, the delivery time shown at checkout is the first promise a customer
receives, and it is made before any courier has accepted the order. If the promise is too short, orders
arrive late and customers lose trust; if it is padded, customers wait longer than necessary or order
elsewhere. Behind the promise sits an operational question as well: knowing which orders will take
long, and why, tells the operations team where and when extra couriers are needed. Large platforms
build such estimates from rich internal data~\cite{zhu2020fulfillment}, and evaluate them on that same data,
so it is hard to tell how far order-time information alone can go.

We study this with the public Meituan dataset released for the INFORMS TSL data-driven research
challenge~\cite{zhao2024meituan} (eight days, 23 districts). Two choices define the problem. First, we predict the \emph{delivery time} (minutes from
order placement to arrival), not the delay relative to the platform's promise; delay is used only to
describe the data and to score promises. Second, the prediction is made \emph{at order placement}:
anything recorded later, such as the assigned courier, dispatch and pickup times, is never a model
input. The platform's promise, though shown at checkout, is also kept out, so that our estimate stays
independent of it and can be compared with it.

\medskip
\noindent\fbox{\parbox{0.96\columnwidth}{\textbf{Research question.} At the moment an order is placed,
using only information known at that moment, how accurately can we predict its delivery time, and which
order-time conditions drive long deliveries?}}
\medskip

We set two success criteria before modelling. The technical one is a test-set mean absolute error (MAE)
at least 20\% below a median predictor, with a bootstrap interval that excludes zero. The business one
is that an 80th-percentile model, used as a promise, is late less often than the platform when both
promise the same average time.

\textbf{Contributions.}
\begin{itemize}
  \item A leakage-safe, time-split comparison of three model families against two baselines: test MAE
        31\% below the median predictor, level with the platform's de-padded estimate.
  \item Real-time features (rider availability, restaurant history and three others) that lower the
        cross-validated MAE by a further 0.33 minutes.
  \item A ranking of delay drivers, an analysis of where the model fails, and a leakage audit of the
        information that does not exist at ordering.
  \item A promise test at equal mean promised time: better estimates alone do not give fewer late
        orders.
\end{itemize}
